\documentclass{amsart}
\input{../aux/style}
% !TEX root = ../axiomatic.tex

% Font encoding and typography
\usepackage[T1]{fontenc}
\usepackage{lmodern}
\usepackage{microtype}

% Mathematics
\usepackage{mathtools}
\usepackage{amssymb}
\usepackage{mathbbol} % Extended \mathbb{} support

% Colors
\usepackage[dvipsnames]{xcolor}

% Diagrams
\usepackage{tikz-cd}

% Text and quoting
\usepackage{csquotes}
\usepackage[shortcuts]{extdash} % Hyphenation, e.g. non\-/constructively

% Lists
\usepackage{enumerate}
\usepackage{enumitem}

% Comments and to-do's
\usepackage{todo}
\def\todoname{Comments, remarks and to-do's}

\input{../aux/usualcmds}
\input{../aux/commands}
\addbibresource{../aux/usualpapers.bib}
\addbibresource{../aux/bibliography.bib}

\title[Additional material]{THE SPACE OF CUP-$i$ CONSTRUCTIONS \\ \large Additional material for ``An axiomatic characterization of Steenrod's cup-$i$ products''}

\author{Anibal~M.~Medina-Mardones}

\begin{document}
\begin{abstract}
	This companion note collects material complementing the main text.
	We prove that the independence of the freeness axiom persists when the non-zero axiom is strengthened to non-degeneracy, constructing an explicit non-degenerate, irreducible, non-free simplicial \mbox{cup-$i$} construction from a mod~$2$ Pascal transform.
	We also record first results and computer experiments on the structure of the space of all irreducible semi-simplicial \mbox{cup-$i$} constructions.
\end{abstract}
\maketitle

\section{Introduction}\label{s:intro}

Throughout, \emph{the main text} refers to the paper ``An axiomatic characterization of Steenrod's \mbox{cup-$i$} products,'' whose notation, conventions, and results we use freely; \cref{s:recollections} recalls what is needed.

The main text proves that a semi-simplicial \mbox{cup-$i$} construction that is non-zero, irreducible, and free is isomorphic to the canonical one, and that the three axioms are independent.
This note explores the landscape around that statement in two directions.
First, we consider the natural strengthening of the non-zero axiom to \emph{non-degeneracy} and show that the freeness axiom remains independent in this more demanding axiom system (\cref{t:non_deg_non_free}); the counterexample witnessing this is built from a mod~$2$ Pascal transform and is, to our knowledge, the first explicit non-degenerate irreducible \mbox{cup-$i$} construction not isomorphic to the canonical one.
Second, we begin a study of the space of all irreducible semi-simplicial \mbox{cup-$i$} constructions, classifying those supported in a single dimension (\cref{p:level_supported}) and recording computer experiments suggesting a complete description (\cref{s:space}).

\section{Recollections}\label{s:recollections}

As shown in the main text, a semi-simplicial \mbox{cup-$i$} construction $\triangle$ is equivalent to a collection of elements $\triangle_i[n] \in \cP(\simplex^n)^{\ot 2}_{i+n}$, indexed by $i, n \in \N$, satisfying
\begin{equation}\label{eq:defining}
	\bd\, \triangle_i[n] + \triangle_i \bd\,[n] = (1+T)\, \triangle_{i-1}[n],
\end{equation}
where $\triangle_i\bd\,[n] = \sum_{j=0}^{n} \cP(\delta_j)^{\ot 2}\, \triangle_i[n-1]$ and $\triangle_{-1} = 0$.
Since these identities are linear, the set of semi-simplicial \mbox{cup-$i$} constructions is an $\F$-vector space, closed under $\triangle \mapsto T\triangle$.
For a chain $\zeta$ and a basis element $V \ot W$ of $\cP(\simplex^n)^{\ot 2}$ we write $\coeff{\zeta}{V \ot W} \in \F$ for the coefficient of $V \ot W$ in $\zeta$.
A basis element $V \ot W$ is \emph{irreducible} if $V \cap W = \emptyset$ and \emph{reducible} otherwise, with $\pired$ denoting the projection onto the span of the reducible ones; a construction is \emph{irreducible} if all summands of all $\triangle_i[n]$ are, and \emph{free} if no $\triangle_i[n]$ with $i \neq n$ contains both a basis element and its transposition among its summands.
The \emph{canonical} construction is $\canonical_i[n] = \sum_{U \in \P^n_{n-i}} U^0 \ot U^1$, where $U = U^0 \sqcup U^1$ is the splitting by the index function.

For $U \subseteq \set{0, \dots, n}$, the \emph{partition chain} $\zeta_U$ is the sum of all basis elements $V \ot W$ with $V \union W = U$ and $V \cap W = \emptyset$.
We import the following results from the main text, where they are established in the complex associated to the augmented simplex category; the degree hypothesis below accounts for the passage to the ordinary complex.

\begin{lemma}[Partition chains]\label{l:partition_chains}
	A homogeneous irreducible chain $\zeta \in \cP(\simplex^n)^{\ot 2}$ of degree at least $n$ satisfies $(\pired \circ \bd)\zeta = 0$ if and only if $\zeta$ is a sum of partition chains.
	Furthermore, $\bd\, \zeta_U = \sum_{\bar u \notin U} \zeta_{\bar u . U}$, with the summand indexed by the full subset interpreted as its truncation to non-negative bidegrees.
\end{lemma}

\begin{theorem}[Simpliciality]\label{t:irreducible_implies_simplicial}
	Every irreducible semi-simplicial \mbox{cup-$i$} construction is simplicial.
\end{theorem}

We will also use the coefficient comparison extracted from the proof of the simpliciality theorem:
for a collection $\set{\zeta[n]}_n$ of chains $\zeta[n] \in \cP(\simplex^n)^{\ot 2}$ supported on irreducible basis elements, the identity $\pired\big(\bd\, \zeta[n]\big) = \sum_j \cP(\delta_j)^{\ot 2}\, \zeta[n-1]$ holds if and only if
\begin{equation}\label{eq:coefficient_comparison}
	\coeff{\zeta[n]}{X \setminus k \ot Y} + \coeff{\zeta[n]}{X \ot Y \setminus k} =
	\coeff{\zeta[n-1]}{\sigma_k(X \setminus k) \ot \sigma_k(Y \setminus k)}
\end{equation}
for every basis element $X \ot Y$ with $X \cap Y = \set{k}$, where $\sigma_k$ denotes the order-preserving relabeling of $\set{0, \dots, n} \setminus \set{k}$ onto $\set{0, \dots, n-1}$, extended elementwise to subsets:
indeed, the reducible summands of $\bd\, \zeta[n]$ and all summands of the coface sum are supported on basis elements $X \ot Y$ whose factors intersect in a single element $k$, arising respectively from the two lifts $X {\setminus} k \ot Y$ and $X \ot Y {\setminus} k$, and from the unique preimage under $\cP(\delta_k)^{\ot 2}$.

\section{Non-degeneracy}\label{s:nondeg}

A natural strengthening of the non-zero axiom is \defn{non-degeneracy}: $x \cup_0 x \neq 0$ for every $0$-simplex $x$.
Since a non-degenerate construction is in particular non-zero, the main theorem holds verbatim with non-degenerate in place of non-zero: there is up to isomorphism only one non-degenerate, irreducible and free semi-simplicial \mbox{cup-$i$} construction.

The independence of the axioms, however, does not transfer automatically, since a counterexample must satisfy all the remaining axioms, including the strengthened one.
Two of the three examples of the main text serve both axiom systems: the zero construction is degenerate, and the transposing example is non-degenerate since it agrees with $\canonical$ on $0$-simplices.
In contrast, $(1+T)\canonical_n[n] = (1+T)(\emptyset \ot \emptyset) = 0$ shows that $(1+T)\canonical$ satisfies $x \cup_0 x = 0$ for every $0$-simplex $x$, so it is excluded by non-degeneracy, and the independence of freeness requires a new example.
Constructing one is the goal of the next two sections.

\section{Columns}\label{s:columns}

\begin{lemma}\label{l:columns}
	Fix $i_0 \in \N$ and let $\Theta = \set[\big]{\Theta[n]}_{n\in\N}$ be a collection of elements $\Theta[n] \in \cP(\simplex^n)^{\ot 2}_{n + i_0}$ satisfying, for every $n \in \N$,
	\[
	(1+T)\,\Theta[n] = 0
	\qquad \text{and} \qquad
	\bd\, \Theta[n] + \Theta\bd\,[n] = 0,
	\]
	where $\Theta\bd\,[n] = \sum_{j=0}^{n} \cP(\delta_j)^{\ot 2}\, \Theta[n-1]$.
	Then the collection $\triangle$ defined by $\triangle_{i_0}[n] = \Theta[n]$ and $\triangle_{i}[n] = 0$ for $i \neq i_0$ is a semi-simplicial \mbox{cup-$i$} construction, referred to as the \defn{column} of $\Theta$ at $i_0$.
\end{lemma}

\begin{proof}
	The defining identity \eqref{eq:defining} at $(i_0, n)$ reduces to the second hypothesis, the one at $(i_0+1, n)$ to the first, and all others hold trivially.
\end{proof}

For example, for each $i_0 \in \N$ the collection $\set[\big]{(1+T)\canonical_{i_0}[n]}_{n}$ satisfies both hypotheses, and $(1+T)\canonical$ is the sum of the associated columns.

\section{A non-degenerate non-free construction}\label{s:pascal}

In this section we construct a column $\Theta$ at $0$ such that $\canonical + \Theta$, provided by \cref{l:columns} and linearity, is non-degenerate, irreducible, and not free.

For functions $b \colon \N^2 \to \F$ define the \defn{Pascal transform}
\[
(\rB b)(w,v) = \sum_{w' \leq w} \sum_{v' \leq v} \binom{w}{w'} \binom{v}{v'}\, b(w',v').
\]
The transform $\rB$ is an involution: since $\binom{w}{w'}\binom{w'}{w''} = \binom{w}{w''}\binom{w - w''}{w' - w''}$, summing over $w'$ gives $\sum_{w'} \binom{w}{w'}\binom{w'}{w''} = \binom{w}{w''}\, 2^{w - w''} \equiv \delta_{w w''} \bmod 2$, and similarly in the second variable.

Let $\rho \colon \N^2 \to \F$ be the indicator function of $\set{(1,1),\, (0,2)}$ and define
\[
c = \rho + \rB\rho,
\]
so that $\rB c = \rB \rho + \rho = c$.
Explicitly, $(\rB\rho)(w,v) = wv + \binom{v}{2}$, so that
\[
c(w,v) =
\begin{cases}
	wv + \binom{v}{2} \bmod 2 & \text{if } w + v \geq 3, \\
	\hfil 0 & \text{if } w + v \leq 2.
\end{cases}
\]

For $n \in \N$ and $u \in \set{0, \dots, n}$ write $\hat{u} = \set{0, \dots, n} \setminus \set{u}$ and, for $V \subseteq \hat u$, define
\[
\theta_n(u; V) =
\sum_{A \subseteq V} c\Big(u - \bars[\big]{A \cap \set{0, \dots, u-1}},\ n - u - \bars[\big]{A \cap \set{u+1, \dots, n}}\Big) \in \F.
\]
Since every basis element $V \ot W$ of $\cP(\simplex^n)^{\ot 2}_{n}$ with $V \cap W = \emptyset$ satisfies $V \union W = \hat u$ for a unique $u$, the element
\[
\Theta[n] =
\sum_{u = 0}^{n} \ \sum_{\substack{V \union W = \hat{u} \\ V \cap W = \emptyset}} \theta_n(u; V)\ V \ot W \ \in \ \cP(\simplex^n)^{\ot 2}_{n}
\]
is a sum of irreducible basis elements.

Let us record the first values of $\Theta$.
If $n \leq 2$, every summand of $\theta_n(u;V)$ has $c$-arguments $(w,v)$ with $w + v \leq n \leq 2$, so $\Theta[n] = 0$.
If $n = 3$, the summands with $A \neq \emptyset$ vanish for the same reason and, since $c(0,3) = c(1,2) = 1$ and $c(2,1) = c(3,0) = 0$, we obtain $\theta_3(u;V) = c(u, 3-u) = 1$ exactly when $u \in \set{0,1}$, independently of $V$.
Therefore
\begin{equation}\label{eq:theta_3}
	\Theta[3] = \zeta_{\set{1,2,3}} + \zeta_{\set{0,2,3}}.
\end{equation}

\begin{lemma}\label{l:theta_derivative}
	For distinct $u, k \in \set{0, \dots, n}$ and $V \subseteq \hat{u} \setminus \set{k}$,
	\[
	\theta_n(u;\, V \union \set{k}) + \theta_n(u;\, V) =
	\theta_{n-1}\big(\sigma_k(u);\, \sigma_k(V)\big).
	\]
\end{lemma}

\begin{proof}
	In the sum defining $\theta_n(u; V \union \set{k})$, the terms with $k \notin A$ add up to $\theta_n(u; V)$.
	The remaining terms are indexed by $A = A' \union \set{k}$ with $A' \subseteq V$, and their $c$-arguments are obtained from those of $A'$ by decreasing the first coordinate by one if $k < u$, and the second if $k > u$.
	On the other hand, $\sigma_k$ restricts to an order-preserving bijection from $\set{0, \dots, n} \setminus \set{k}$ to $\set{0, \dots, n-1}$, so the assignment $A' \mapsto \sigma_k(A')$ is a bijection onto the subsets of $\sigma_k(V)$ satisfying
	\[
	\bars[\big]{\sigma_k(A') \cap \set{0, \dots, \sigma_k(u) - 1}} = \bars[\big]{A' \cap \set{0, \dots, u - 1}},
	\qquad
	\bars[\big]{\sigma_k(A') \cap \set{\sigma_k(u) + 1, \dots, n-1}} = \bars[\big]{A' \cap \set{u+1, \dots, n}}.
	\]
	Since $\sigma_k(u) = u - 1$ if $k < u$ and $\sigma_k(u) = u$ if $k > u$, and consequently $(n-1) - \sigma_k(u)$ equals $n - u$ if $k < u$ and $n - u - 1$ if $k > u$, the $c$-arguments of $\sigma_k(A')$ in $\theta_{n-1}\big(\sigma_k(u); \sigma_k(V)\big)$ coincide with those of $A' \union \set{k}$ above, and the claim follows.
\end{proof}

\begin{theorem}\label{t:non_deg_non_free}
	The collection $\Theta$ satisfies the hypotheses of \cref{l:columns}, and $\canonical + \Theta$, with $\Theta$ regarded as a column at $0$, is a simplicial \mbox{cup-$i$} construction that is non-degenerate, irreducible, and not free.
	In particular, the freeness axiom cannot be removed from the main theorem after strengthening non-zero to non-degeneracy.
\end{theorem}

\begin{proof}
	\textit{Symmetry.}
	We show by induction on $n$ that $g_n(u; V) := \theta_n(u; V) + \theta_n(u; \hat u \setminus V)$ vanishes identically; this gives $(1+T)\Theta[n] = 0$.
	For $n \leq 2$ this holds since $\Theta[n] = 0$.
	Applying \cref{l:theta_derivative} to both summands of $g_n$,
	\[
	g_n(u; V \union \set{k}) + g_n(u; V) = g_{n-1}\big(\sigma_k(u); \sigma_k(V)\big),
	\]
	which vanishes by the induction hypothesis, so $g_n(u; -)$ is a constant function.
	Its value at $V = \emptyset$ is
	\[
	\theta_n(u; \emptyset) + \theta_n(u; \hat u) =
	c(u, n-u) + (\rB c)(u, n-u) = 0,
	\]
	using $\rB c = c$ for the last equality.

	\medskip
	\textit{Chain identity, reducible part.}
	Since the summands of $\Theta[n]$ are irreducible, the identity $\pired\big(\bd\, \Theta[n]\big) = \Theta\bd\,[n]$ is equivalent to \cref{eq:coefficient_comparison} holding for $\set{\Theta[n]}_n$.
	Writing $u$ for the unique element of $\set{0,\dots,n} \setminus (X \union Y)$, that identity reads
	\[
	\theta_n(u;\, X \setminus k) + \theta_n(u;\, X) =
	\theta_{n-1}\big(\sigma_k(u);\, \sigma_k(X \setminus k)\big),
	\]
	which is \cref{l:theta_derivative} applied to $V = X \setminus \set{k}$.

	\medskip
	\textit{Chain identity, irreducible part.}
	Since $\Theta\bd\,[n]$ is a sum of reducible basis elements, it remains to show that $\pi^\perp_{\text{red}}\big(\bd\,\Theta[n]\big) = 0$.
	The irreducible summands of $\bd\, \Theta[n]$ arise by adjoining the missing vertex $u$ to one of the two factors, so they are supported on basis elements $P \ot Q$ with $P \union Q = \set{0, \dots, n}$ and $P \cap Q = \emptyset$, and the coefficient of such an element is
	\[
	\Phi_n(P, Q) :=
	\sum_{u \in P} \theta_n(u;\, P \setminus u) \ + \
	\sum_{u \in Q} \theta_n(u;\, P),
	\]
	an expression meaningful for every partition $P \sqcup Q = \set{0,\dots,n}$, including the improper ones with $P$ or $Q$ empty.
	We show by induction on $n$ that $\Phi_n$ vanishes identically; for $n \leq 2$ this holds since $\Theta$ does.
	For $k \notin P \union Q$ with $P \sqcup Q = \set{0,\dots,n} \setminus \set{k}$, the terms $\theta_n(k; P)$ appearing in both $\Phi_n(P \union \set{k},\, Q)$ and $\Phi_n(P,\, Q \union \set{k})$ cancel, and \cref{l:theta_derivative} applied to the remaining terms gives
	\[
	\Phi_n(P \union \set{k},\, Q) + \Phi_n(P,\, Q \union \set{k}) =
	\Phi_{n-1}\big(\sigma_k(P),\, \sigma_k(Q)\big) = 0,
	\]
	so $\Phi_n$ is constant, and it suffices to evaluate it at $(P, Q) = \big(\set{0, \dots, n-1},\, \set{n}\big)$.
	The term with $u = n$ is $\theta_n(n; P) = \sum_{A \subseteq P} c(n - \bars{A}, 0) = 0$, since $c(w, 0) = 0$ for every $w$.
	For $u < n$, counting the subsets $A$ by $j = \bars{A \cap \set{0,\dots,u-1}}$ and $l = \bars{A \cap \set{u+1, \dots, n-1}}$,
	\[
	\Phi_n(P,Q) = \sum_{u=0}^{n-1} \sum_{j \leq u} \sum_{\ \ \mathclap{l \leq n-1-u}} \binom{u}{j} \binom{n-1-u}{l}\, c(u - j,\ n - u - l).
	\]
	We treat the three summands of the decomposition $c(w,v) = \rho(w,v) + wv + \tbinom{v}{2}$ separately, assuming $n \geq 4$; the case $n = 3$ follows directly from $\theta_3(u;V) = c(u, 3-u)$ and $c(0,3) + c(1,2) + c(2,1) + c(3,0) = 0$.
	The $\rho$-summand contributes only through the arguments $(u-j,\, n-u-l) = (1,1)$, requiring $(j, l) = (u-1,\, n-u-1)$ with coefficient $\binom{u}{u-1} = u$, and $(0,2)$, requiring $(j,l) = (u,\, n-u-2)$ with coefficient $\binom{n-1-u}{n-u-2} = n-1-u$; its total contribution is $\sum_{u=0}^{n-1} \big(u + (n - 1 - u)\big) = n(n-1) \equiv 0$.
	For the $\tbinom{v}{2}$-summand, the inner sum over $j$ contributes $\sum_j \binom{u}{j} = 2^u$, reducing the total to the single value $u = 0$, where it equals
	\[
	\sum_{l} \binom{n-1}{l}\binom{n-l}{2} =
	\big[z^{n-2}\big]\, (1+z)^{n-1} (1+z)^{-3} =
	\big[z^{n-2}\big]\, (1+z)^{n-4} = 0 \bmod 2,
	\]
	using $\sum_{s} \binom{s+2}{2} z^s = (1+z)^{-3}$ over $\F$ and $n - 2 > n - 4 \geq 0$.
	For the $wv$-summand, the sum factors as
	\[
	\sum_{u=0}^{n-1} \Big(\sum_{j} \binom{u}{j}(u-j)\Big) \Big(\sum_{l} \binom{n-1-u}{l}(n-u-l)\Big),
	\]
	whose first factor equals $u\, 2^{u-1}$ for $u \geq 1$ and $0$ for $u = 0$, hence is odd only for $u = 1$, and whose second factor, with $M = n - 1 - u$, equals $(M+2)\,2^{M-1}$ for $M \geq 1$ and $1$ for $M = 0$, hence is odd only for $u \in \set{n-2, n-1}$; since $n \geq 4$ no $u$ satisfies both, and the $wv$-summand vanishes.
	Therefore $\Phi_n \equiv 0$, completing the proof that $\Theta$ is a column at $0$.

	\medskip
	\textit{Conclusion.}
	By linearity and \cref{l:columns}, $\widetilde\triangle := \canonical + \Theta$ is a semi-simplicial \mbox{cup-$i$} construction.
	It is irreducible since both summands are, and hence simplicial by \cref{t:irreducible_implies_simplicial}.
	It is non-degenerate since $\Theta[0] = 0$.
	It is not free: by \cref{eq:theta_3}, both $\set{1} \ot \set{2,3}$ and $\set{2,3} \ot \set{1}$ are summands of $\Theta[3]$, while the only summand of $\canonical_0[3]$ supported on partitions of $\set{1,2,3}$ is $\set{1,2,3} \ot \emptyset$, so both remain summands of $\widetilde\triangle_0[3]$.
	Finally, since freeness is preserved by isomorphisms of \mbox{cup-$i$} constructions -- an isomorphism replaces each $\triangle_i$ by either $\triangle_i$ or $T\triangle_i$ -- the construction $\widetilde\triangle$ is not isomorphic to the canonical one.
\end{proof}

\section{The quotient by reducible chains}\label{s:quotient}

The main text identifies $\widehat{\cP}(\simplex^n)$ with the Koszul complex $(\Lambda_n, \omega \cdot -)$ of $\omega = x_0 + \dots + x_n$ in the exterior algebra on the vertices, the partition chains with the image of its coproduct $\Delta$, and the span of the reducible basis elements with the kernel of its multiplication $\mu$, an ideal and hence a subcomplex.
We record three complementary observations.

First, the fibers of $\mu$ over the monomial basis are the partitions: $\mu(x_V \ot x_W) = x_U$ precisely when $V \ot W$ is a partition of $U$.
The partition coproduct is therefore the transfer of $\mu$, summing each fiber, and the composite
\[
\mu\, \Delta(x_U) = 2^{\bars{U}}\, x_U
\]
vanishes for $U \neq \emptyset$, the mod~$2$ shadow of the usual transfer identity.

Second, consider the quotient $\rQ(\simplex^n)$ of $\cP(\simplex^n)^{\ot 2}$ by the reducible subcomplex, which we identify with the irreducible subspace equipped with the differential $\piirred \circ \bd$.
For an irreducible semi-simplicial \mbox{cup-$i$} construction, every summand of the coface term $\triangle_i\bd\,[n]$ is reducible, whereas $(1+T)\triangle_{i-1}[n]$ is irreducible, so passing to the quotient the defining identity becomes
\[
\bd_{\rQ}\, \big[\triangle_i[n]\big] = (1+T) \big[\triangle_{i-1}[n]\big].
\]
An irreducible construction therefore determines, for each $n$, a $\Sym_2$-equivariant chain map from $W$ to $\rQ(\simplex^n)$.

Third, computer calculations for $n \leq 4$ show that the homology of $\rQ(\simplex^n)$ is one-dimensional, concentrated in degree $n$, and generated by the class of the Alexander--Whitney diagonal $\canonical_0[n]$.
The following sketch explains this topologically.
Under the equivalence with chains, the reducible subcomplex corresponds to $\sum_F \chains(F)^{\ot 2}$, the sum over the facets $F$ of $\simplex^n$, a model for the union of the products $F \times F$ inside $\simplex^n \times \simplex^n$; the nerve of this cover of the union is the boundary of an $n$-simplex, its pieces and their intersections are contractible, so the union has the homotopy type of a sphere of dimension $n-1$, and the long exact sequence of the quotient gives $\rH_k\big(\rQ(\simplex^n)\big) \cong \widetilde{\rH}_{k-1}(S^{n-1})$.
Moreover, reading the removed sets as vertex sets, the irreducible basis elements are exactly the cells of the simplicial \emph{deleted product} of $\simplex^n$ -- the pairs of vertex-disjoint faces -- a classical object of embedding theory carrying a free $\Sym_2$-action away from the diagonal, where equivariant cohomology and the resolution $W$ are the standard tools.
In this light, an irreducible \mbox{cup-$i$} construction is an equivariant chain-level representative of the generator of $\rH_n\big(\rQ(\simplex^n)\big)$, and the main theorem describes the rigidity of such representatives in the presence of the freeness axiom.

\subsection*{A Hopf-algebraic program}

As recalled in the main text, the algebra $\Lambda_n$ is the group ring $\F[C_2^{\, n+1}]$ under $t_u = 1 + x_u$, but the partition coproduct is the primitively generated structure, the Cartier dual of the group-ring one.
Several threads deserve a fuller treatment, which we leave for future work.
\begin{enumerate}
	\item The primitively generated exterior Hopf algebra is self-dual: the top class $x_{\set{0,\dots,n}}$ furnishes a Poincar\'e duality pairing $\coeff{x_U \cdot x_{U'}}{x_{\set{0,\dots,n}}}$ under which the multiplication $\mu$ and the partition coproduct $\Delta$ are adjoint.
	The duality between the reducible subcomplex $\ker\mu$ and the partition subcomplex $\operatorname{im}\Delta$, and hence between the description of the partition chains in \cref{l:partition_chains} and the reducible quotient, is the algebraic form of the Poincar\'e--Lefschetz duality of the deleted product sketched above.
	\item The complex $W$ is the minimal free resolution of $\F$ over $\F[C_2]$, and $\widehat{\cP}(\simplex^n)$ is the Koszul complex of the augmentation element over $\F[C_2^{\, n+1}] = \F[C_2]^{\ot(n+1)}$.
	Both sides of the defining identity are thus multiplication by an augmentation element in the group ring of an elementary abelian $2$-group, and a \mbox{cup-$i$} construction is an equivariant map between such complexes.
	We caution that $\widehat{\cP}(\simplex^n)$ is the \emph{finite} Koszul complex, not the infinite periodic resolution $W^{\ot(n+1)}$ of $\F$ over the same ring, nor is the relevant coproduct the group-ring one; conflating these would misidentify the structure.
	\item The odd-primary theory would replace $C_2$ by $C_p$ and the exterior algebra by a truncated polynomial algebra, connecting the present picture to the cup-$(p,i)$ constructions of \cite{medina2021may_st}.
\end{enumerate}

\section{The space of cup-\textit{i} constructions}\label{s:space}

The examples of the main text and the construction of \cref{s:pascal} indicate that the space of semi-simplicial \mbox{cup-$i$} constructions is large.
Since the defining identities are linear, the collection of irreducible semi-simplicial \mbox{cup-$i$} constructions with $\triangle_i[n]$ prescribed only for $n \leq N$ -- referred to as \emph{truncated} at $N$ -- forms an $\F$-vector space, and restriction defines linear maps from each truncation level to the previous one.
The following result classifies the elements supported in a single level, which constitute the kernels of these restriction maps.

\begin{proposition}\label{p:level_supported}
	Let $\set{X_i}_{0 \leq i \leq n}$, $X_i \in \cP(\simplex^n)^{\ot 2}_{n+i}$, be irreducible chains and consider the truncated construction with $\triangle_i[n] = X_i$ and $\triangle_i[m] = 0$ for $m < n$.
	It satisfies the defining identities if and only if every $X_i$ is a sum of partition chains and, under the correspondence $\zeta_U \mapsto U$, the resulting collection is a cycle in the augmented chain complex of $\simplex^n$.
	The space of such collections has dimension $2^n - 1$.
\end{proposition}

\begin{proof}
	Since all lower levels vanish, the coface terms of the defining identities vanish, and the identities reduce to $\bd X_i = (1+T) X_{i-1}$ for $0 \leq i \leq n+1$, with $X_{-1} = X_{n+1} = 0$.
	The right hand sides are sums of irreducible basis elements, so $(\pired \circ \bd) X_i = 0$ and, by \cref{l:partition_chains}, each $X_i$ is a sum of partition chains.
	Partition chains are fixed by $T$, so $(1+T)X_{i-1} = 0$, and the identities reduce further to $\bd X_i = 0$ for every $i$.
	Under $\zeta_U \mapsto U$, the boundary formula of \cref{l:partition_chains} identifies the complex of sums of partition chains with the chain complex of $\simplex^n$ augmented by the class of $\zeta_{\set{0, \dots, n}}$ in degree $-1$: for $\bars{U} = n$ the only element not in $U$ produces the partition chain of the full set, so the condition $\bd X_0 = 0$ includes the augmentation.
	Therefore, the collections in question correspond to the cycles of the augmented chain complex $\widetilde{\chains}(\simplex^n)$.
	This complex is acyclic, so its cycles agree with its boundaries in every degree, and writing $\widetilde{Z}_i$ for the space of degree-$i$ cycles we obtain $\dim \widetilde{\chains}_i = \dim \widetilde{Z}_i + \dim \widetilde{Z}_{i-1}$ for $0 \leq i \leq n$.
	Summing over $i$ and using $\sum_{i \geq 0} \dim \widetilde{\chains}_i = 2^{n+1} - 1$, $\dim \widetilde{Z}_{-1} = 1$, and $\dim \widetilde{Z}_n = 0$ gives $\sum_{i \geq 0} \dim \widetilde{Z}_i = 2^n - 1$.
\end{proof}

Computer experiments complement this result; the scripts, self-contained and independent of any library, are available in the repository of this project.
The dimension of the space of irreducible constructions truncated at $N$ equals
\[
2^{N+1} - N - 1
\qquad \text{for } N \leq 6,
\]
that is, $5, 12, 27, 58, 121$ for $N = 2, \dots, 6$, and every restriction map in this range is surjective: every truncated irreducible construction extends to the next level.
Note that these dimensions are consistent with \cref{p:level_supported}: the increment from $N-1$ to $N$ equals $2^N - 1$.
Moreover, imposing the codegeneracy conditions changes none of these dimensions, in agreement with \cref{t:irreducible_implies_simplicial}.

These observations suggest the following problems, which we leave open:
\begin{enumerate}
	\item Prove that every truncated irreducible construction extends to the next level; together with \cref{p:level_supported}, this would establish the dimension formula for all $N$ and show that the space of irreducible semi-simplicial \mbox{cup-$i$} constructions is infinite dimensional.
	\item Classify the columns at a fixed index $i_0$, that is, the natural symmetric irreducible cycle families of \cref{l:columns}.
	Computer experiments indicate that, at $i_0 = 0$, beyond the twist family $(1+T)\canonical_0$ there is a space of solutions whose first non-vanishing components are sums $\zeta_{\hat{u}} + \zeta_{\hat{v}}$ of partition chains of co-vertex sets, of which the construction of \cref{s:pascal} is one member; a closed-form description of all of them is unknown.
	\item Describe the quotient of the space of irreducible constructions by the isomorphism relation, that is, by the twist action of $\prod_{\N} \Sym_2$.
\end{enumerate}

\printbibliography
\end{document}
